\documentclass[10pt,a4paper]{amsart}
\usepackage[margin=19mm]{geometry}
\usepackage{amsmath,amssymb,mathtools}
\usepackage[scr=rsfs,cal=euler]{mathalfa}
\usepackage{xfrac}
\usepackage[unicode,hidelinks,bookmarks=false]{hyperref}
\newcommand{\ie}{i.e.\ }
\newcommand{\cf}{cf.\ }
\newcommand{\resp}{respectively\ }
\newcommand{\Subsection}{\subsection}
\providecommand{\nobreakdash}{\nobreak\hbox{-}}
\setlength{\emergencystretch}{2em}
\multlinegap0pt
\usepackage{bm}
\usepackage{bbm}
\usepackage{dsfont}
\usepackage{xspace}
\usepackage{cancel}
\usepackage{mathtools}
\usepackage{amsthm}
\makeatletter
\@ifundefined{theorem}{\newtheorem{theorem}{Theorem}}{}
\makeatother
\title[In Search of the Canonical Harmony for 12-TET]{In Search of the Canonical Harmony for 12-TET}
\author{Pawel Nurowski}
\date{Source: arXiv:2601.02271v3 (CC BY 4.0)}
\begin{document}
\maketitle

\noindent\emph{Source and licence.} In Search of the Canonical Harmony for 12-TET. Version arXiv:2601.02271v3 (CC BY 4.0), distributed under the Creative Commons Attribution 4.0 International licence, as recorded in the arXiv licence metadata for this exact version and shown on the abstract page https://arxiv.org/abs/2601.02271v3. Full rights notice: licenses/arxiv-2601.02271-CC-BY-4.0.txt.\par\smallskip

\noindent\textit{Editorial scope.} Counted unit: the Theorem whose source label is \texttt{None} in the article (source0.tex lines 1217--1219, source label \texttt{None}), delivered together with the article's own complete original proof of it (source0.tex lines 1221--1242, one \texttt{proof} environment of 22 lines). No mathematics is written, repaired or re-derived: statement and proof are the article's own text line for line. Required context retained uncounted: none. Every definition, display and label of those lines is retained unchanged. Omitted from this package and not redistributed: 1--1216, 1220--1220, 1243--5045. Nothing inside a retained excerpt is omitted or altered: every retained body is delivered exactly as the article's lines are written, so no omission receipt of any kind accompanies this package. Citations: the retained excerpts carry no citation command at all, so no bibliography entry of the article is transcribed and no reference is redistributed. Reference adaptation: none was needed and none was made; every reference inside the delivered unit resolves inside it, so no unresolved reference or citation is produced. Printed slips kept literal: no printed word, symbol, hypothesis or number of the counted statement or of its proof was altered. Layout and class: the article's own class is \texttt{amsart} with packages including fontenc, inputenc, babel, lmodern, placeins, float, geometry, hyperref, booktabs, enumitem, url, xcolor, colortbl, amsmath; the wrapper is the lane's pinned \texttt{amsart} editorial wrapper at 10pt on A4, which restates the theorem-like environments the retained text needs and adds the article's own macro definitions that the retained text needs (none), with extra wrapper packages (bm, bbm, dsfont, xspace, cancel, mathtools). The delivered numbering is the unit's own. Assets: no retained span contains a figure environment, a graphics-inclusion command, a PDF-page inclusion, a table or a TikZ picture; no image file is copied into this package, the package declares no asset dependency and no image file is redistributed with it. Licence: version arXiv:2601.02271v3 of this article is distributed under the Creative Commons Attribution 4.0 International licence, as recorded in the arXiv licence metadata for this exact version and shown on the abstract page https://arxiv.org/abs/2601.02271v3. A full rights notice is delivered at licenses/arxiv-2601.02271-CC-BY-4.0.txt. Characters: every retained excerpt is pure ASCII, so the delivered engine, XeLaTeX with its default Latin Modern text font, reports no missing character.
\begin{theorem}
In the restored cubic system, there are exactly two isomorphism classes of minimal cycles (of length 4) under the rotational symmetry of \texorpdfstring{$\mathbb{Z}_{12}$}{Z12}. These classes are distinguished by their chirality (direction of winding through the wormhole).
\end{theorem}

\begin{proof}
A cycle must alternate between Major and Minor chords. The minimal length is 4. A 4-cycle must involve at least two wormhole crossings (\texorpdfstring{$L$}{L}) to leave and return to the starting universe.
The two minimal solutions correspond to the two directions of traversing the wormhole:

\begin{enumerate}
    \item \textbf{The Violet Cycle (Right-handed):} $ S = L_{+5} \cdot P \cdot L_{+7} \cdot R_{+2} $.\\
    Tracing the index change starting from $M_0$: 
    $M_0 \xrightarrow{L} m_5 \xrightarrow{P} M_5 \xrightarrow{L} m_{10} \xrightarrow{R} M_0$.
    The total displacement is $+5 + 0 + 5 + 2 = +12 \equiv 0 \pmod{12}$.
    
    \item \textbf{The Orange Cycle (Left-handed):} $ S' = R_{-2} \cdot L_{+7} \cdot P \cdot L_{+5} $. \\
    We trace the cycle $M_4 \to m_2 \to M_9 \to m_9 \to M_4$:
    \begin{itemize}
        \item $M_4 \xrightarrow{R} m_2$ (displacement $-2$)
        \item $m_2 \xrightarrow{L} M_9$ (displacement $+7$)
        \item $M_9 \xrightarrow{P} m_9$ (displacement $0$)
        \item $m_9 \xrightarrow{L} M_4$ (displacement $+7$. Note: $9+7=16 \equiv 4$)
    \end{itemize}
    Total displacement: $-2 + 7 + 0 + 7 = +12 \equiv 0 \pmod{12}$.
\end{enumerate}
These two sequences are non-isomorphic because they traverse the graph with opposite chirality.
\end{proof}

\end{document}
